\chapter{Limitations of the Work}

\vspace*{0.3cm}

The project provides a useful optimizer comparison, but the results must be
interpreted within the limits of the experimental setup.

\section{Dataset and Training Duration}

The processed dataset is intentionally small: 2000 training examples, 500
validation examples, and 500 test examples. This size is suitable for a light
academic experiment, but a larger dataset would provide more reliable conclusions.

\vspace*{0.15cm}

The benchmark was also short, with 10 training steps and a batch size of 4. It
compares early optimizer behavior rather than complete model convergence. Longer
training would be necessary to confirm the final ranking of the optimizers.

\section{Hardware Dependency}

The benchmark was executed on a machine equipped with a performant GPU. Metrics
such as training latency, throughput, inference time, and VRAM usage depend on
the hardware. The trends are useful, but the exact values may change on another
machine.

\section{Energy and Compression Limits}

Energy consumption is important for Edge AI, but accurate measurement requires
power telemetry. When such telemetry is unavailable, energy analysis remains
limited. In addition, the current work focuses on fine-tuning and optimizer
comparison. Compression methods such as quantization, pruning, and knowledge
distillation are left for future work.

\section{Second-Order Optimizers}

L-BFGS and Newton-CG are included for theoretical comparison, but their high
memory and computational costs make them difficult to use for transformer
fine-tuning. Their results should therefore be interpreted mainly from an
academic optimization perspective.

\section{Conclusion}

These limitations define the scope of the project. The work gives a controlled
comparison of optimizer behavior, while future work should include longer
training, larger datasets, real energy measurements, and compression experiments.
